\documentclass[11pt]{article}
\usepackage[notes]{yalatex}
% \addbibresource{references.bib}   % uncomment for citations

\title{MATH730: Algebraic Topology}
\author{Yuval Amit}
\date{Instructor: Or Hershkovitz \\ Fall 2026}

\begin{document}
\maketitle

\tableofcontents
\newpage

\section{Motivation and Review}

\begin{notebox}[Motivation]
    Distinguish topological objects, and ultimately classify them.
\end{notebox}

Let's see some examples where the power of algebraic topology helps us prove important topological properties.

\begin{thm}
    The unit ball $\bar{\mb D}^{n+1}$ is not homeomorphic to $S^n$
\end{thm}

\begin{proof}
    The $n^{\text{th}}$ homology group $H_n$ distinguishes them. Namely, $H_n\paren{\bar{\mb D}^{n+1}} = \{0\}$, while $H_n(S^n) = \Z$.
\end{proof}

\begin{thm}[][label=thm:boundaryfixedthm]
    There does not exist a continuous map $f:\bar{\mb D}^{n+1} \to S^n$ such that $f(x)=x$ for every $x \in S^n$ (note that $S^n$ is the boundary of $\bar{\mb D}^{n+1}$).
\end{thm}

\begin{proof}
    Suppose such an $f$ exists. Consider the inclusion map $i: S^n \to \bar \D^{n+1}$. Then, the map $f \circ i : S^n \to S^n $ is the identity map. Moreover, $i$ and $f$ induce maps on the homology groups 
    $$H_n(S^n) \arr[H_n(i)] H_n(\D^{n+1}) \arr[H_n(f)] H_n(S^n).$$
    By functoriality,
    $$H_n(f) \circ H_n(i) = H_n(f \circ i) = H_n(\id) = \id.$$
    Hence, we have 
    \begin{align*}
        \Z \to0 \to \Z
    \end{align*}
    induces the identity map $\Z \to \Z$, which is impossible.
\end{proof}

\begin{thm}[Brauer's Fixed Point Theorem]
    Every continuous function $f:\bar{\D}^n \to \bar{\D}^n$ has a point $x$ such that $f(x) = x$.
\end{thm}

\begin{proof}
    Suppose $f$ does not have a fixed point. Then, for every $x \in \bar{\D}^n$, we may consider the unique straight line from $f(x)$ to $x$. Let $g(x)$ be the point where the line $f(x) \to x$ hits $S^n$. Observe that $g(x) = x$ for every $x \in S^n$, and $g$ is continuous. Thus, we have a contradiction by \Cref{thm:boundaryfixedthm}.
\end{proof}

\subsection{Topology Review}

\begin{defn}
    In $\R^n$, we have the usual metric $d(x,y)$. We denote $B(x,r)$ to be the open ball around $x$ with radius $r$.
\end{defn}

\begin{defn}
    A topological space $X$ is called \textbf{compact} if every open cover has a finite subcover.
\end{defn}

\begin{thm}
   A subset $X \subset \R^n$ is compact if and only if $X$ is closed and bounded. 
\end{thm}

\begin{proof}
    Suppose $X$ is compact. The set $\{B(0,n)\}_{n=1}^\infty$ of open balls is necessarily an open cover of $X$, hence has a finite subcover, thus $X$ is bounded. Moreover, if $\R^n-X$ is not open, let $p \in \R^n\setminus X$ be such that $B(p,\veps) \cap X \neq \emptyset$. Then, take the cover of $X$ by 
    $$\brax{B(x,\frac{d(p,x)}{2})}_{x \in X}.$$
    One sees that this does not have a finite subcover, a contradiction.

    For the other direction, note that closed+bounded implies sequential compactness. Let $\{U_\alpha\}_{\alpha \in I}$ be an open cover of $X$. 
\begin{claim}
    There exists $\delta>0$ such that for every $x \in X$, $B(x,\delta) \subset U_\alpha$ for some $\alpha \in I$.
\end{claim}

\begin{proof}
    Contrarily, if this were not the case, there exists a sequence $\{x_i\} \subset X$ such that $B(x,1/n) \not \subset U_\alpha$ for all $\alpha \in I$. Then, $\{x_i\}$ has a subsequence converging to some $p$. This $p$ must be contained in some $U$ in the open cover, and it follows that there exists a $\delta$ and $N$ with $B(x_N,\delta) \subset U$.
\end{proof}

Since $X$ is bounded in $\R^n$, for every $\delta>0$ the number of disjoint balls of radius $\delta$ around points in $X$ is bounded. Let $m$ be this upper bound. Take a point $p_1 \in X$. The ball $B(p,\delta)$ is contained in some cover element $U_1$. If $U_1$ covers $X$, we're done, but otherwise repeat with some $p_2$ and $U_2$. Continue to repeat until we have $p_{m+1}$ and $U_{m+1}$. We claim that $U_1,\dots,U_{m+1}$ must cover $X$: indeed, if it doesn't, then the balls $\{B(p_i,\delta/2)\}_{i=1}^{m+1}$ must all be disjoint, and this is a contradiciton.
\end{proof}

\begin{defn}
    A topological space $X$ is called \textbf{connected} if $X = U_1 \sqcup U_2$ with $U_1,U_2$ open implies $U_1 = \emptyset$ or $U_2 =\emptyset$.
\end{defn}

\begin{defn}
    $X$ is called \textbf{path connected} if for every $p,q \in U$, there exists a continuous map $\delta:[0,1] \to X$ such that $\delta(0)=p$ and $\delta(1) = q$.
\end{defn}

\begin{prop}
    Path connected $\implies$ connected.
\end{prop}

\begin{defn}
    A topological space $X$ is called \textbf{Hausdorff} if for every distinct $p,q \in X$, there exists $p \in U$ and $q \in V$ open such that $U \cap V = \emptyset$.
\end{defn}

\begin{thm}
    Let $f: X \to Y$ be bijective and continuous, with $X$ compact, $Y$ Hausdorff. Then, $f$ is a homeomorphism.
\end{thm}

\begin{proof}
    It suffices to show $f$ is an closed map. Let $C \subset X$ be closed. Since $X$ is compact, one can see that $C$ is compact (theorem). Therefore, $f(C)$ is compact (theorem), hence closed (theorem).
\end{proof}

\begin{prop}
    If $X$ is Hausdorff and $A \subset X$ is compact, then $A$ is closed.
\end{prop}

\begin{exercise}
    Give an example of a space $X$ such that $A \subset X$ is compact but not closed. This space cannot be Hausdorff.
\end{exercise}

\begin{defn}
    Let $X$ be a topological space, and $\sim$ be an equivalence relation on $X$. The \textbf{quotient-topology} is defined as follows: let $\pi: X \to X/\!\!\sim$ be the projection map. Then, $U \subset X/\!\!\sim$ is open whenever $\pi^{-1}(U)$ is open.
\end{defn}

\begin{prop}
    A map $f:(X/\!\!\sim) \to Y$ is continuous if and only if $f \circ \pi$ is continuous.
\end{prop}

\begin{ex}
    Let $\R\P^{n}$ be the space of lines through the origin in $\R^{n+1}$: this is the quotient space of $\R^{n+1}\setminus \{0\}$ by the equivalence $x \sim y \iff$ there exists $\lambda \in \R$ with $\lambda x=y$. This is homeomorphic to the quotient space of $S^n$ by the relation $x \sim -x$.
\end{ex}

\begin{defn}
    Let $(X_\alpha,\mc U_\alpha)_{\alpha \in I}$ be topological spaces. The \textbf{product topology} on $\prod_{I}X_\alpha$ has a basis of sets of the form
    $$\prod_{\alpha \in U}U_\alpha,  \quad U_\alpha \in \mc U_\alpha,$$
    and $U_\alpha =X_\alpha$ for all but finitely many $\alpha \in I$.
\end{defn}

\begin{thm}[Tychonoff's Theorem][label=thm:tychonoff]
    The product of compact spaces is compact.
\end{thm}

\section{CW Complexes}

\begin{defn}
    A \textbf{CW complex} is a topological space $X$, together with a sequence of subspaces $X_0 \subset X_1 \subset \cdots $ such that
    \begin{enumerate}
        \item $X = \bigcup X_n$
        \item $U\subset X$ is open if and only if $U \cap X_n$ is open for all $n \geq 0$
        \item $X_0$ is a discrete set of points
        \item For every $n \geq 1$, there exists an indexing set $I(n)$, disks $\{D^n_\alpha\}_{\alpha \in I(n)}$, and continuous maps
        $\varphi^n_\alpha : \partial D^n_\alpha 
         \to X_{n-1}$ such that 
         $$X_n = \bigsqcup_{\alpha \in I(n)}\bar D_\alpha^n\sqcup X_n  /\sim$$
         where $x \sim \varphi_\alpha^n(x)$ for all $\alpha \in I(n)$, $x \in \partial D_\alpha^n$.
    \end{enumerate}
    We call $X_n$ the $n^{\text{th}}$ \textbf{skeleton} of $X$, and $\varphi_\alpha^n$ is called an \textbf{attaching map}. For a CW complex $X$, the sets $\{X_n\}$ and $\{ \varphi_\alpha^n\}$ are called a \textbf{CW structure} on $X$. The quotient maps $\Phi_\alpha^n:\bar D_\alpha^n \to X_n$ are called the \textbf{characteristic maps}. If $X_{n-1} \subsetneq X_n$, but $X_i=X_j$ for all $i, j\geq n$, we say that $X$ is $\bm n$\textbf{-dimensional}.
\end{defn}


\begin{ex}
    A discrete set of points is a $0$-dimensional CW complex.
\end{ex}

\begin{ex}
    A graph is a $1$-dimensional CW complex.
% https://q.uiver.app/#q=WzAsNyxbMCwxLCJYXzAiXSxbMiwxLCJcXGJ1bGxldCJdLFs0LDEsIlxcYnVsbGV0Il0sWzMsMCwiWF8wIl0sWzMsMywiXFxidWxsZXQiXSxbMiwyLCJcXGJ1bGxldCJdLFsxLDMsIlxcYnVsbGV0Il0sWzAsMCwiIiwwLHsic3R5bGUiOnsiaGVhZCI6eyJuYW1lIjoibm9uZSJ9fX1dLFsxLDIsIiIsMCx7InN0eWxlIjp7ImhlYWQiOnsibmFtZSI6Im5vbmUifX19XSxbMywxLCIiLDAseyJzdHlsZSI6eyJib2R5Ijp7Im5hbWUiOiJkYXNoZWQifX19XSxbMywyLCIiLDIseyJzdHlsZSI6eyJib2R5Ijp7Im5hbWUiOiJkYXNoZWQifX19XSxbNSw0LCIiLDIseyJzdHlsZSI6eyJoZWFkIjp7Im5hbWUiOiJub25lIn19fV0sWzYsNiwiIiwyLHsiYW5nbGUiOi0xMzUsInN0eWxlIjp7ImhlYWQiOnsibmFtZSI6Im5vbmUifX19XSxbNSw2LCIiLDAseyJzdHlsZSI6eyJoZWFkIjp7Im5hbWUiOiJub25lIn19fV0sWzYsNCwiIiwwLHsic3R5bGUiOnsiaGVhZCI6eyJuYW1lIjoibm9uZSJ9fX1dXQ==
\[\begin{tikzcd}[ampersand replacement=\&,cramped]
	\&\&\& {X_0} \& \\
	{X_0} \&\& \bullet \&\& \bullet \\
	\&\& \bullet \\
	\& \bullet \&\& \bullet
	\arrow[dashed, from=1-4, to=2-3]
	\arrow[dashed, from=1-4, to=2-5]
	\arrow[no head, from=2-1, to=2-1, loop, in=55, out=125, distance=10mm]
	\arrow[no head, from=2-3, to=2-5]
	\arrow[no head, from=3-3, to=4-2]
	\arrow[no head, from=3-3, to=4-4]
	\arrow[no head, from=4-2, to=4-2, loop, in=190, out=260, distance=10mm]
	\arrow[no head, from=4-2, to=4-4]
\end{tikzcd}\] 
\end{ex}

\begin{ex}
    $\R^n$ is a CW complex. For example, $\R^2$ has the $0$-cell $\Z^2$, the $1$-cell a $\Z^2$-grid, and the $2$-cell induced by attaching each disk to the square of the grid.
\end{ex}

\begin{ex}
    A torus is a CW complex.
\end{ex}

\begin{ex}
    $S^n$ is a CW complex. Namely, $X_0$ is a single point, $X_1=\cdots = X_{n-1}=X_0$, and the (only possible) attaching map $\varphi_1^n : D_1^n \to X_{n-1}$.
\end{ex}

\begin{ex}
    $\R\P^n=S^n/(x \sim -x)$ is a CW complex. When $n=0$, $\R\P^n$ is just a point, hence as a CW structure. Suppose $\R\P^{n-1}$ has a CW structure. Let $\bar D_\alpha^n$ be the upper hemisphere. We obtain the CW structure on $\R \P^n$ by attaching one $n$-cell to $\R\P^{n-1}$ via the attaching map 
    $$\varphi_\alpha^n(x_1,\dots,x_n,0) = [(x_1,\dots,x_n)] \subset \R\P^{n-1}.$$
\end{ex}

\begin{exercise}
    Find a CW structure on $\C \P^n = (\C^{n+1}\setminus \{0\})/(x \sim \lambda x,\lambda \in \C^*)$
\end{exercise}

\begin{lem}
    Let $X$ be a CW-complex. Then, $U \subset X$ is open $\iff$ for every $n$ and $\alpha \in I(n)$, $(\Phi^n_\alpha)^{-1}(U)$ is open.
\end{lem}

\begin{proof}
    Recall that $U$ is open in $X$ if and only if $U \cap X^n$ is open for all $n$, if and only if for the projection map
    $$P^n: X^n \sqcup \bigsqcup _{\alpha \in I(n)}D^\alpha_n \to X^n,$$
    we have
    $$(P^n)^{-1}(U \cap X^n)$$
    is open. This is the case if and only if for each $\alpha \in I(n)$,
    $$(\Phi_\alpha^n)^{-1}(U \cap X)$$
    and $U \cap X^{n-1}$ are open. So, $(\implies)$ is verified, and $(\impliedby)$ follows by induction on $n$.
\end{proof}

\begin{lem}
    Let $K\subset X$ be compact. Then, $K$ intersects only finitely many cells in $X$. 
\end{lem}

\begin{proof}
    Suppose not. Then, we can choose an infinite set $S \subset K$ such that no two points in $S$ belong to the same cell. We claim $S$ is closed. Indeed, note that $S \cap X^0$ is closed since $X^0$ is discrete. By induction, suppose $S \cap X^{n-1}$ is closed. Then, $S \cap X^n$ is closed $\iff$ $(\Phi_\alpha^n)^{-1}(S)$ is closed, and by induction $(\Phi_\alpha^n)^{-1}(S) \cap \partial D_\alpha^n$ is closed, so indeed $(\Phi_\alpha^n)^{-1}(S)$ is closed as the union of closed sets. Thus, $S \cap X^n$ is closed for every $n$, hence $S$ is closed. 
    
    In fact, we've shown that every subset of $S$ is closed, so $S$ is discrete. Moreover, $S$ is compact as a closed subset of a compact set. Thus, $S$ is compact and discrete, hence finite.
\end{proof}

\begin{defn}
    Let $A \subset X$ be a set, and fix $\veps :=\{\veps_\alpha^n \in \R^+\}$ for every $n$ and $\alpha \in I(n)$. Then, an $\bm \veps$\textbf{-neighborhood} of $A$ is the open set $N_\veps(A) \subset X$ defined as follows: let
    \begin{align*}
        N_\veps^0(A) := A \cap X^0
    \end{align*}
    and then for every $\alpha \in I(n)$, set
    $$(\Phi_\alpha^n)^{-1}(N_\veps^n(A)) := \braks{B(A \cap \Int(D_\alpha^n),\veps_\alpha^n) \cap \Int(D_\alpha ^n)} \cup \braks{ (1-\veps _\alpha^n,1](\Phi_\alpha^n)^{-1}\paren{N_\veps^{n-1}(A)}},$$
    and
    $$N_\veps(A) = \bigcup N_\veps^n(A).$$
\end{defn}

\begin{prop}
    $$N_\veps(A) \cap X^n = N_\veps^n(A)$$
\end{prop}

\begin{thm}
    Let $X$ be a CW complex. Then, $X$ is Hausdorff.
\end{thm}

\begin{proof}
    Let $p \neq q \in X$, and suppose that $p,q \in X^n$ but $q \not \in X^{n-1}$. So, $q$ is in some disk $e_\alpha^n$. If $p \in X^{n-1}$, we may choose $\veps$ sufficiently small so that $N_\veps(X^{n-1})$ and $N_\veps(q)$ are disjoint. If $p \not \in X^{n-1}$, then $p,q$ are distinct points in $e_\alpha^n$ which is Hausdorff, hence separable in $\veps$-neighborhoods. These $\veps$-neighborhoods lift to the global structure of $X$ so that $N_\veps(p) \cap N_\veps(q) = \emptyset$.
\end{proof}

\begin{defn}
    Let $X,Y$ be topological spaces. Two (continuous) maps $f,g:X \to Y$ are \textbf{homotopic} if there exists $H:X \times [0,1] \to Y$ continuous such that $H(x,0)=f(x)$ and $H(x,1)=g(x)$ for every $x \in X$.

    For a subset $A \subset X$, two maps $f,g:X \to Y$ are called \textbf{homotopic relative to} $\bm A$ if $f|_A=g|_A$, and there exists a homotopy $H:X \times [0,1]\to Y$ between $f,g$ such that $H(x,t)=f(x)=g(x)$ for all $x \in A$, $t \in [0,1]$.
\end{defn}
Homotopy is denoted $f \sim g$ or $f \overset{H}{\sim} g$.

\begin{defn}
    Two maps $f:X \to Y$ and $g:Y \to X$ are called \textbf{homotopy equivalent} if $g \circ f \sim \id_X$ and $f \circ g = \id_Y$. In this case, $X$ and $Y$ are said to be \textbf{homotopically equivalent.}

    A subset $A \subset X$ is called a \textbf{retract} of $X$ if there exists $r:X \to A$ continuous such that $r(x)=x$ for every $x \in A$. Moreover, $A$ is called a \textbf{deformation retract} if $r \sim \id_X$ relative to $A$. $X$ is called \textbf{contractible} if $\id_X \sim $ constant
\end{defn}

\begin{ex}
    Show that $\R^{n+1}\setminus \{0\}$ deformation retracts to $S^n \subset \R^{n+1}$.
\end{ex}

\begin{ex}
    Show that $i:\R^n \to \R^{n+k}$ and $\pi:\R^{n+k}\to \R^n$ are homotopy equivalent.
\end{ex}

\begin{ex}
    Let $X$ be a topological space. Define $C(X) = X \times [0,1] / [(x,1) \sim (y,1)]$. Show that $C(X)$ is contractible.
\end{ex}

\begin{defn}
    $X$ is called \textbf{locally contractible} if for every $p \in X$ and $U \ni p$ open, there exists $V \subset U$ with $p \in V$ such that $p$ is contractible.
\end{defn}

\begin{prop}
    Every CW complex is locally contractible.
\end{prop}

\begin{proof}
    Exercise. Inductively push the additional part of the neighborhood to the boundary.
\end{proof}

\begin{cor}
    If $X$ is a CW complex, then $X$ is connected $\iff$ it's path connected.
\end{cor}

\begin{proof}
    Exercise. Local contractibility allows you to understand path connectivity because of the path induced by a retract.
\end{proof}

\section{Fundamental Group}

\begin{defn}
    Let $X$ be a topological space $p \in X$. The \textbf{fundamental group} of $X$ at $p$, is
    $$\pi_1(X,p):=\brax{ \gamma:[0,1]\to X \mid \gamma(0)=\gamma(1) } / \braks{ \gamma \overset{\text{rel}(\{0,1\})}{\sim}\mu }.$$
    Equivalently, let $q$ be a fixed point in $S^1$, and we have
    $$\pi_1(X,p) = \brax{\gamma :S^1 \to X}/\braks{\gamma \overset{\text{rel(q)}}{\sim} \mu}$$
    This definition is in some ways better, since changing $S^1$ to $S^n$ gives a more generalized definition of the $\bm n$\textbf{-th fundamental group} $\pi_n(X,p)$.
\end{defn}

\begin{defn}
    A connected space $X$ is called \textbf{simply connected} if there exists $p \in X$ such that $\pi_1(X,p)=\{ [p]\}$.
\end{defn}

\begin{prop}
    If $X$ is contractible, then
    $$\pi_1(X,p) =\{[p]\}$$
    for the $p$ such that $X$ deformation retracts to $p$.
\end{prop}

\begin{proof}
    Let $\gamma \in \pi_1(X,p)$. Let $H$ be a homotopy between $\id_X$ and $p$. Take
    $$\widetilde{H}(s,t) := H(\gamma(s),t).$$
    Then, $\gamma \overset{\widetilde H}{\sim}p$.
\end{proof}

\begin{defn}
    For $\gamma,\mu \in \pi_1(X,p)$, define
    $$\gamma \star \mu(t) = \case{\mu(2t) & 0 \leq t \leq 1/2 \\ \gamma (2t-1) & 1/2 \leq t \leq 1}.$$
    Then, $\star$ induces a well-defined group structure on $\pi_1(X,p)$.
\end{defn}

Indeed, in order to show that this concatenation of curves defines a group structure on $\pi_1(X,p)$, one must prove the following claim:

\begin{claim}
    If $[\gamma] = [\gamma']$ and $[\mu] = [\mu ']$, then $[\gamma \star \mu] = [\gamma ' \star \mu ']$. Thus, $[\gamma] \star [\mu]$ is a well-defined binary operation on $\pi_1(X,p)$. Moreover, $\pi_1(X,p)$ is a group, with identity the constant map $[p]$, and $[\gamma(t)]^{-1}=[\gamma(1-t)]$.
\end{claim}

\begin{exercise}
    Show that if $\varphi :[0,1] \to [0,1]$ is a continuous map with $\varphi(0)=0$ and $\varphi(1)=1$, then $\gamma \circ \varphi \sim \gamma$.
\end{exercise}

\begin{thm}[][label=thm:pathconnectedisom]
    Let $X$ be path connected, $p,q\in X$. Then, $\pi_1(X,p) \cong \pi_1(X,q)$.
\end{thm}

\begin{proof}
Since $X$ is path connected, we have a continuous curve $\mu:[0,1] \to X$ with $\mu(0)=p$ and $\mu(1)=q$.
    \begin{claim}
        Let $\Omega_p$ denote the space of curves starting and ending at $p$. Then, the map
    $$T_\mu :\Omega_p \to \Omega_q$$
    defined by $T_\mu := \mu \star \gamma \star \mu^{-1}$
    induces an isomorphism $\pi_1 (X,p) \arr[\sim] \pi_1(X,q)$.
    \end{claim}
    Indeed, suppose $[\gamma] = [\gamma ']$. Then, $[T_\mu (\gamma)] = [T_\mu(\gamma')]$. Moreover,
    $$\mu \star (\gamma _1 \star \gamma_2)\star \mu^{-1}=\mu \star \gamma _1 \star \mu^{-1}\star \mu \star \gamma_2 \star \mu^{-1},$$
    hence $T_\mu$ is a homomorphism. Moreover, the map $\psi =\mu(1-t)$ induces a clear inverse homomorphism $T_{\psi}:\pi_1(X,q) \to \pi_1(X,p)$, so it follows that $T_\mu$ is an isomorphism.
\end{proof}

\begin{prop}
    Let $(X,p)$ and $(Y,q)$ be (pointed) topological spaces, and $f: (X,p) \to (Y,q)$ be continuous (i.e. a continuous map $f:X \to Y$ with $f(p)=q$). Then, the map $f_*:=\pi_1(f):\pi_1(X,p) \to \pi_1(Y,q)$ defined by
    $$f_*([\gamma]) = [f \circ \gamma]$$
    is a well-defined group homomorphism. Moreover, if
    $$(X,p) \arr[f] (Y,q) \arr[g] (Z,r)$$
    is a chain of continuous maps, then $(g\circ f)_* = g_* \circ f_*$.
\end{prop}

\begin{proof}
    Let $[\gamma] = [\gamma ']$. Then, $f\circ \gamma \sim f\circ \gamma'$ on $f \circ H$, relative to $\{0,1\}$, so this is well-defined, and
    $$(f \circ \gamma ) \star (f \circ \mu) = f(\gamma \star \mu)$$
    at the curves level, which induces the homomorphism
    $(g \circ f)_* =g_* \circ f_*$
    at the level of fundamental groups.
\end{proof}

\begin{thm}[Functoriality of the Fundamental Group]
The map $\pi_1$ on spaces and the continuous maps between them is a functor from the category of pointed spaces to the category of groups. That is,
$$\pi_1 : \w{Top_\bullet} \to \w{Grp}$$
is a covariant functor.
\end{thm}

\begin{defn}
    $X$ is called \textbf{strongly contractable} if there exists $p \in X$ such that $X$ deformation retracts to $p$.
\end{defn}

\begin{prop}
    If $X$ is strongly contractable, then $\pi_1(X,q) = \{e\}$ for all $q \in X$.
\end{prop}

\begin{proof}
When $q=p$, we know $\id _X \overset{\w{rel} p}{\sim}(x \mapsto p)$, and it follows that $\pi_1(X,p)$ is trivial. Then, the theorem follows for arbitrary $q$ by \Cref{thm:pathconnectedisom}, since strongly contractable implies path connected.
\end{proof}

\begin{notebox*}[Recall]
    $X$ is called contractible if $X \sim \{q\}$, i.e. $\id_X \sim (x \maps q)$.
\end{notebox*}

\begin{thm}
    If $X$ is contractible, then $\pi_1(X,q) = \{e\}$ for all $q \in X$.
\end{thm}

\section{\texorpdfstring{The Fundamental Group of $S^1$}{The Fundamental Group of S1}}

\begin{thm}[][label=thm:circlelift]
    Let $\gamma:[0,1] \to S^1$ be continuous with $\gamma(0)=1$. Then, there exists a unique $\theta:[0,1] \to \R$ such that $\theta (0)=0$, and
    $$e^{i\theta(t)}=\gamma(t).$$
    Moreover, we can do this continuously along homotopies. That is, if $H:[0,1]\times [0,1] \to S^1$ has $H(0,s)=1$, there there exists a unique $\theta: [0,1] \times [0,1] \to \R$ such that $\theta(0,s)=0$ for all $s$, and $e^{i\theta}=H$ for all $t,s$.
\end{thm}


\begin{cor}
    Let $\gamma_1,\gamma_2:[0,1] \to S^1$ be such that $\gamma_i(0)=\gamma_i(1)=1$, $i=1,2$. Let $\theta_1,\theta_2$ be their lifts as in \Cref{thm:circlelift}. Then,
    $$\gamma_1 \overset{\w{rel}\{0,1\}}{\sim}\gamma_2 \iff \theta_1(1)=\theta_2(1)$$
\end{cor}

\begin{proof}
($\implies$) Suppose $\gamma_1 \sim \gamma_2$ relative to $\{0,1\}$, via a homotopy $H$. Let $\theta:[0,1]\times [0,1] \to \R$ be the lift of $H$ as in \Cref{thm:circlelift}. Then, $\theta_1(x)=\theta(x,0)$ and $\theta_2(x)=\theta(x,1)$. Consider the path $s \mapsto \theta(1,s)$. Then, since $H(1,s)=1$, and $e^{i\theta}=H$, we have $\theta(1,s) \in 2 \pi \Z$. So, as $\theta(1,s)$ is continuous with a discrete image, it must be constant, hence $\theta_1(1)=\theta_2(1)$.

$(\impliedby)$ Take $\theta(t,s)= (1-s)\theta_1(t) + s\theta_2(t)$. Then, $\theta(t,s)$ is a homotopy between $\theta_1$ and $\theta_2$, with $\theta(1,s)=\theta_1(1)$, and $\theta(0,s)=0$. Now, we may take $H(t,s):= e^{i\theta(t,s)}$, and $H$ is the desired homotopy from $\gamma_1$ to $\gamma_2$ relative to $\{0,1\}$. Note that relative to $1$ comes from the fact $\theta_1(1)=\theta_2(1)$.
\end{proof}

\begin{cor}
For any $p \in S^1$,
    $$\pi_1(S^1,p) \cong \Z,$$
    with $t \mapsto e^{2\pi i t}:[0,1]\to S^1$ being the generator.
\end{cor}

\begin{proof}
    Clearly, $(e^{2\pi i t})^k \overset{\w{rel}\{0,1\}}{\sim} e^{2\pi i kt}$, where the LHS denotes concatenation $k$ times. On the other hand, each element in $\pi_1(S^1,1)$ is in bijective correspondence with $t \mapsto e^{2\pi i kt}$ for some $k \in \Z$. Indeed, if $\theta$ is the lift of an element in $\pi_1(S^1,1)$, then $\theta_1(1) \in 2\pi \Z$, and the bijective correspondence is induced by $\gamma \leftrightarrow \frac{\theta_1(1)}{2\pi}$.
\end{proof}

\begin{exercise}
    Let $\{ (X_i,p_i)\}_{i=1}^N$ be a set of pointed spaces. Show that
    $$\pi_1\paren{\prod_{k=1}^N X_i, (p_1,\dots,p_n)} \cong \prod_{k=1}^N \pi_1(X_i,p_i).$$
\end{exercise}

\begin{exercise}
    Show that a continuous map $f:\bar D^2 \to S^1$ satisfying $f(x)=x$ for all $x \in S^1$ does not exist.
\end{exercise}

\begin{prop}
    For all $k \geq 2$, $\pi_1(S^k,p)$ is trivial.
\end{prop}

\begin{proof}
    Let $\gamma:[0,1] \to S^k$ be continuous. We claim $\gamma$ is homotopic, relative to $\{0,1\}$, to a curve $\tilde \gamma$ that misses a point $q$. The proposition follows from the claim: indeed, we have $\psi:S^k \setminus \{q\} \arr[\sim] \D^k $, so consider $\mu:=\pi \circ \tilde \gamma$, which is homotopic to the constant map, and we may push it back to $S^k$ via $\psi^{-1}$ to show $[\gamma]=[\tilde \gamma]=[\id]$.

    To prove the claim, we approximate $\gamma$ via smooth maps, and smooth maps always miss points by \href{https://en.wikipedia.org/wiki/Sard%27s_theorem}{Sard's theorem}.
\end{proof}


\section{Amalgamated Products and Free Products}
\begin{defn}
    Let $H$, $G_1$, $G_2$ be groups, and let $\varphi_i:H \to G_i$ be homomorphisms. The \textbf{amalgamated product} of $G_1$ and $G_2$ along $H$, denoted $G_1*_HG_2$, is a group together with homomorphisms $\psi_i:G_i \to G_1 *_HG_2$ satisfying $\psi_1 \circ \varphi_1=\psi_2 \circ \varphi_2$. Moreover, it has the universal property that for every group $K$ and homomorphisms $f_i:G_i\to K$ with $f_1 \circ \varphi_1 = f_2 \circ \varphi_2$, there exists a unique homomorphism $T:G_1*_HG_2 \to k$ such that $T \circ \psi_i=f_i$. That is, the following diagram commutes:
    % https://q.uiver.app/#q=WzAsNSxbMCwwLCJIIl0sWzIsMCwiR18xIl0sWzAsMiwiR18yIl0sWzIsMiwiR18xKl9IR18yIl0sWzMsMywiSyJdLFswLDEsIlxcdmFycGhpXzEiXSxbMCwyLCJcXHZhcnBoaV8yIiwyXSxbMSwzLCJcXHBzaV8xIiwyXSxbMiwzLCJcXHBzaV8yIl0sWzEsNCwiZl8xIl0sWzIsNCwiZl8yIiwyXSxbMyw0LCJcXGV4aXN0cyAhVCIsMSx7InN0eWxlIjp7ImJvZHkiOnsibmFtZSI6ImRhc2hlZCJ9fX1dXQ==
\[\begin{tikzcd}[ampersand replacement=\&,cramped]
	H \&\& {G_1} \& \\
	\\
	{G_2} \&\& {G_1*_HG_2} \\
	\&\&\& K
	\arrow["{\varphi_1}", from=1-1, to=1-3]
	\arrow["{\varphi_2}"', from=1-1, to=3-1]
	\arrow["{\psi_1}"', from=1-3, to=3-3]
	\arrow["{f_1}", from=1-3, to=4-4]
	\arrow["{\psi_2}", from=3-1, to=3-3]
	\arrow["{f_2}"', from=3-1, to=4-4]
	\arrow["{\exists !T}"{description}, dashed, from=3-3, to=4-4]
\end{tikzcd}\]
\end{defn}

A good way to think about this is gluing $G_1$ and $G_2$ along $H$.

\begin{ex}
    Consider a similar definition but with sets $A \arr[\varphi_1]B$ and $A \arr[\varphi_2] B_2$. If we have $f_1:B_1 \to K$ and $f_2:B_2 \to K$, then one can see $B_1 \sqcup B_2$ with the natural embeddings $i_1:B_1 \to B_1 \sqcup B_2$, $i_2:B_2\to B_1 \sqcup B_2$, has the desired unique $\tilde T:B_1 \sqcup B_2 \to K$. However, we do not have the desired commutativity $i_1 \circ \varphi_1 = i_2 \circ \varphi_2$. In order to have this, we need to quotient $B_1 \sqcup B_2$ by the relation $i_1\varphi_1(x) \sim i_2\varphi_2(x)$. Then,
    $(B_1 \sqcup B_2)/\sim$ is the amalgamated product of these sets.
\end{ex}

\begin{prop}[Existence and Uniqueness of the Amalgamated Product]
    Let $H,G_1,G_2$ be groups with homomorphisms $\varphi_i:H \to G_i$. Then, the amalgamated product $(G_1*_HG_2, \psi_1,\psi_2)$ exists, and if $(\tilde{G_1*_HG_2},\tilde\psi_1,\tilde\psi_2)$ is another amalgamated product, then there exists a unique isomorphism
    $$T:G_1 *_HG_2 \arr \tilde{G_1 *_HG_2}$$
    such that $T \circ \psi_i = \tilde \psi_i$.
\end{prop}

\begin{proof}
    % https://q.uiver.app/#q=WzAsNSxbMCwwLCJIIl0sWzIsMCwiR18xIl0sWzAsMiwiR18yIl0sWzIsMiwiR18xKl9IR18yIl0sWzQsNCwiXFx0aWxkZXtHXzEqX0hHXzJ9Il0sWzAsMiwiXFx2YXJwaGlfMiJdLFswLDEsIlxcdmFycGhpXzEiLDJdLFsxLDMsIlxccHNpXzEiLDJdLFsyLDMsIlxccHNpXzIiXSxbMSw0LCJcXHRpbGRle1xccHNpXzF9Il0sWzIsNCwiXFx0aWxkZXtcXHBzaV8yfSIsMl0sWzMsNCwiXFxleGlzdHMgIVQiLDEseyJzdHlsZSI6eyJ0YWlsIjp7Im5hbWUiOiJhcnJvd2hlYWQifSwiYm9keSI6eyJuYW1lIjoiZGFzaGVkIn19fV1d
\[\begin{tikzcd}[ampersand replacement=\&,cramped]
	H \&\& {G_1} \&\& \\
	\\
	{G_2} \&\& {G_1*_HG_2} \\
	\\
	\&\&\&\& {\tilde{G_1*_HG_2}}
	\arrow["{\varphi_1}"', from=1-1, to=1-3]
	\arrow["{\varphi_2}", from=1-1, to=3-1]
	\arrow["{\psi_1}"', from=1-3, to=3-3]
	\arrow["{\tilde{\psi_1}}", from=1-3, to=5-5]
	\arrow["{\psi_2}", from=3-1, to=3-3]
	\arrow["{\tilde{\psi_2}}"', from=3-1, to=5-5]
	\arrow["{\exists !T}"{description}, dashed, tail reversed, from=3-3, to=5-5]
\end{tikzcd}\]
\end{proof}

\begin{defn}
    Let $G_1,G_2$ be two groups. The \textbf{free product} $G_1*G_2$ is the group of words in $G_1\setminus \{e_1\}$, $G_2\setminus \{e_2\}$, and a new identity $\{e\}$, where elements from the same group multiply as expected, the product of elements is concatenation, and multiplying an element $g$ from $G_1$ or $G_2$ by its inverse in that group yields $e$.
\end{defn}

\begin{thm}
    Let $G_1,G_2$ be groups. Then,
    $$G_1*G_2 \cong G_1*_{\{e\}}G_2$$
\end{thm}

\begin{proof}
The claim follows immediately from the following lemma.
\begin{lem}
    The free product $G_1*G_2$ satisfies the following universal property:
    % https://q.uiver.app/#q=WzAsNSxbMCwwLCJcXHtlXFx9Il0sWzIsMCwiR18xIl0sWzAsMiwiR18yIl0sWzIsMiwiR18xKkdfMiJdLFs0LDQsIksiXSxbMCwyLCIiLDAseyJzdHlsZSI6eyJ0YWlsIjp7Im5hbWUiOiJob29rIiwic2lkZSI6InRvcCJ9fX1dLFswLDEsIiIsMix7InN0eWxlIjp7InRhaWwiOnsibmFtZSI6Imhvb2siLCJzaWRlIjoidG9wIn19fV0sWzEsMywiXFxwc2lfMSIsMix7InN0eWxlIjp7InRhaWwiOnsibmFtZSI6Imhvb2siLCJzaWRlIjoidG9wIn19fV0sWzIsMywiXFxwc2lfMiIsMCx7InN0eWxlIjp7InRhaWwiOnsibmFtZSI6Imhvb2siLCJzaWRlIjoidG9wIn19fV0sWzEsNCwiZl8xIl0sWzIsNCwiZl8yIiwyXSxbMyw0LCJcXGV4aXN0cyAhVCIsMSx7InN0eWxlIjp7ImJvZHkiOnsibmFtZSI6ImRhc2hlZCJ9fX1dXQ==
\[\begin{tikzcd}[ampersand replacement=\&,cramped]
	{\{e\}} \&\& {G_1} \&\& \\
	\\
	{G_2} \&\& {G_1*G_2} \\
	\\
	\&\&\&\& K
	\arrow[hook, from=1-1, to=1-3]
	\arrow[hook, from=1-1, to=3-1]
	\arrow["{\psi_1}"', hook, from=1-3, to=3-3]
	\arrow["{f_1}", from=1-3, to=5-5]
	\arrow["{\psi_2}", hook, from=3-1, to=3-3]
	\arrow["{f_2}"', from=3-1, to=5-5]
	\arrow["{\exists !T}"{description}, dashed, from=3-3, to=5-5]
\end{tikzcd}\]
\end{lem}
\end{proof}

\begin{exercise}
    Show that
    $$G_1*_HG_2 \cong {(G_1*G_2)}/{\angles{\varphi_1(h) \varphi_2(h)^{-1}}}$$
\end{exercise}


\begin{rmk}
    We can similarly define free products $*_\alpha G_\alpha$ of arbitrarily many groups $\{G_\alpha\}_{\alpha \in I}$.
\end{rmk}

\section{Van Kampen's Theorem}

\begin{thm}[Van Kampen's Theorem]
    Let $(X,p)$ be a connected pointed space, and suppose
    $$X = \bigcup_{\alpha \in I}(U_\alpha,p)$$
    where $U_\alpha \subset X$ is open and $p \in U_\alpha$. Suppose that for every $\alpha,\beta \in I$, $U_\alpha \cap U_\beta$ is path connected. Then, the map 
    \begin{align*}
        \Phi: &*_{\alpha}\pi_1(U_\alpha,p) \to \pi_1(X,p)
        \\ &*_{\alpha}[\gamma_\alpha]\mapsto \star_\alpha[\gamma _\alpha]
    \end{align*}
    is surjective. Moreover, if $\alpha,\beta, \delta \in I$ and $U_\alpha \cap U_\beta \cap U_\delta$ is path-connected, then
    $$\pi_1(X,p) \cong *_\alpha\pi_1(U_\alpha, p)/\angles{\angles{ i_{\alpha ,\beta}(\gamma) i_{\beta,\alpha}(\gamma^{-1})}},$$
    where
    $$i_{\alpha,\beta}:\pi_1(U_\alpha \cap U_\beta, p) \to \pi_1(U_\alpha, p)$$
    is the expected embedding, and $\angles{\angles{-}}$ is the normal closure. In particular, if $(X,p) = U_1 \cup U_2$ with $U_1 \cap U_2$ path connected, then
    $$\pi_1(X,p) = \pi_1(U_1,p) \underset{\pi_1(U_1\cap U_2,p)}{*} \pi_1(U_2,p)$$
\end{thm}

\begin{ex}
    Take $X= \bigvee_{\alpha \in I}(S^1,p)$, which is the space defined by attaching $I$ copies of $S^1$ at one shared point $p$. For each $\alpha$, take $U_\alpha = S^1_\alpha \cup N_\veps (p)$. Then,
    $$U_\alpha \cap U_\beta \cap U_\delta = N_\veps(p)$$
    for every $\alpha,\beta,\delta \in I$ is connected (and contractible). Thus, by van Kampen's theorem,
    $$\pi_1(X,p) \cong *_\alpha \pi_1(S^1_\alpha,p)\cong *_\alpha \Z$$
    Indeed, the normal subgroup is trivial because $N_\veps(p)$ is contractible.
\end{ex}

\begin{ex}
    Let $G=(V,E)$ be an undirected connected graph, viewed as a CW-complex. Let $T$ be a spanning tree. Then, deforming the tree to a point, we get a bouquet as in the previous example:
    
    % https://q.uiver.app/#q=WzAsOCxbMSwwLCJcXGJ1bGxldCJdLFswLDEsIlxcYnVsbGV0Il0sWzEsMiwiXFxidWxsZXQiXSxbMiwxLCJcXGJ1bGxldCJdLFszLDIsIlxcYnVsbGV0Il0sWzIsMywiXFxidWxsZXQiXSxbNCwzLCJcXGJ1bGxldCJdLFs2LDEsIlxcYnVsbGV0IixbMTIwLDYwLDYwLDFdXSxbMCwxLCIiLDAseyJjb2xvdXIiOlsxMjAsNjAsNjBdLCJzdHlsZSI6eyJoZWFkIjp7Im5hbWUiOiJub25lIn19fV0sWzEsMiwiIiwwLHsic3R5bGUiOnsiaGVhZCI6eyJuYW1lIjoibm9uZSJ9fX1dLFsyLDMsIiIsMCx7InN0eWxlIjp7ImhlYWQiOnsibmFtZSI6Im5vbmUifX19XSxbMywwLCIiLDAseyJzdHlsZSI6eyJoZWFkIjp7Im5hbWUiOiJub25lIn19fV0sWzEsMSwiIiwwLHsiYW5nbGUiOi05MCwic3R5bGUiOnsiaGVhZCI6eyJuYW1lIjoibm9uZSJ9fX1dLFsyLDQsIiIsMCx7InN0eWxlIjp7ImhlYWQiOnsibmFtZSI6Im5vbmUifX19XSxbMyw0LCIiLDAseyJzdHlsZSI6eyJoZWFkIjp7Im5hbWUiOiJub25lIn19fV0sWzQsNSwiIiwwLHsic3R5bGUiOnsiaGVhZCI6eyJuYW1lIjoibm9uZSJ9fX1dLFs1LDYsIiIsMCx7InN0eWxlIjp7ImhlYWQiOnsibmFtZSI6Im5vbmUifX19XSxbNiw0LCIiLDAseyJzdHlsZSI6eyJoZWFkIjp7Im5hbWUiOiJub25lIn19fV0sWzAsMywiIiwwLHsib2Zmc2V0IjotMywiY29sb3VyIjpbMTIwLDYwLDYwXSwic3R5bGUiOnsiaGVhZCI6eyJuYW1lIjoibm9uZSJ9fX1dLFswLDEsIiIsMCx7Im9mZnNldCI6LTMsInN0eWxlIjp7ImhlYWQiOnsibmFtZSI6Im5vbmUifX19XSxbMywyLCIiLDAseyJvZmZzZXQiOi0zLCJjb2xvdXIiOlsxMjAsNjAsNjBdLCJzdHlsZSI6eyJoZWFkIjp7Im5hbWUiOiJub25lIn19fV0sWzMsNCwiIiwwLHsib2Zmc2V0IjotMywiY29sb3VyIjpbMTIwLDYwLDYwXSwic3R5bGUiOnsiaGVhZCI6eyJuYW1lIjoibm9uZSJ9fX1dLFs0LDUsIiIsMCx7Im9mZnNldCI6LTMsImNvbG91ciI6WzEyMCw2MCw2MF0sInN0eWxlIjp7ImhlYWQiOnsibmFtZSI6Im5vbmUifX19XSxbNCw2LCIiLDAseyJvZmZzZXQiOi0zLCJjb2xvdXIiOlsxMjAsNjAsNjBdLCJzdHlsZSI6eyJoZWFkIjp7Im5hbWUiOiJub25lIn19fV0sWzcsNywiIiwwLHsib2Zmc2V0IjozLCJyYWRpdXMiOi01LCJzdHlsZSI6eyJoZWFkIjp7Im5hbWUiOiJub25lIn19fV0sWzcsNywiIiwwLHsib2Zmc2V0IjozLCJyYWRpdXMiOi01LCJhbmdsZSI6LTkwLCJzdHlsZSI6eyJoZWFkIjp7Im5hbWUiOiJub25lIn19fV0sWzcsNywiIiwwLHsib2Zmc2V0IjozLCJyYWRpdXMiOi01LCJhbmdsZSI6LTE4MCwic3R5bGUiOnsiaGVhZCI6eyJuYW1lIjoibm9uZSJ9fX1dLFs3LDcsIiIsMCx7Im9mZnNldCI6MywicmFkaXVzIjotNSwiYW5nbGUiOjkwLCJzdHlsZSI6eyJoZWFkIjp7Im5hbWUiOiJub25lIn19fV1d
\[\begin{tikzcd}[ampersand replacement=\&,cramped]
	\& \bullet \&\&\&\&\& \\
	\bullet \&\& \bullet \&\&\&\& \textcolor{rgb,255:red,92;green,214;blue,92}{\bullet} \\
	\& \bullet \&\& \bullet \\
	\&\& \bullet \&\& \bullet
	\arrow[color={rgb,255:red,92;green,214;blue,92}, no head, from=1-2, to=2-1]
	\arrow[shift left=3, no head, from=1-2, to=2-1]
	\arrow[shift left=3, color={rgb,255:red,92;green,214;blue,92}, no head, from=1-2, to=2-3]
	\arrow[no head, from=2-1, to=2-1, loop, in=145, out=215, distance=10mm]
	\arrow[no head, from=2-1, to=3-2]
	\arrow[no head, from=2-3, to=1-2]
	\arrow[shift left=3, color={rgb,255:red,92;green,214;blue,92}, no head, from=2-3, to=3-2]
	\arrow[no head, from=2-3, to=3-4]
	\arrow[shift left=3, color={rgb,255:red,92;green,214;blue,92}, no head, from=2-3, to=3-4]
	\arrow[shift right=3, no head, from=2-7, to=2-7, loop, in=310, out=230, distance=15mm]
	\arrow[shift right=3, no head, from=2-7, to=2-7, loop, in=40, out=320, distance=15mm]
	\arrow[shift right=3, no head, from=2-7, to=2-7, loop, in=130, out=50, distance=15mm]
	\arrow[shift right=3, no head, from=2-7, to=2-7, loop, in=220, out=140, distance=15mm]
	\arrow[no head, from=3-2, to=2-3]
	\arrow[no head, from=3-2, to=3-4]
	\arrow[no head, from=3-4, to=4-3]
	\arrow[shift left=3, color={rgb,255:red,92;green,214;blue,92}, no head, from=3-4, to=4-3]
	\arrow[shift left=3, color={rgb,255:red,92;green,214;blue,92}, no head, from=3-4, to=4-5]
	\arrow[no head, from=4-3, to=4-5]
	\arrow[no head, from=4-5, to=3-4]
\end{tikzcd}\]
\end{ex}

\begin{rmk}
    Recall that we required $U_\alpha \cap U_\beta$ to be connected for surjectivity to hold. If they are not connected, then we get examples of maps which aren't surjective: for instance, taking two arclengths of $S^1$, their intersection has two connected components, and $\{e\} *\{e\} \not \surj \Z$.

    The triple intersection condition is also necessary. Consider the graph:
    % https://q.uiver.app/#q=WzAsNSxbMSwwLCJwIl0sWzAsMSwiQSJdLFsxLDEsIkIiXSxbMiwxLCJDIl0sWzEsMiwicSJdLFsxLDQsIiIsMCx7InN0eWxlIjp7ImhlYWQiOnsibmFtZSI6Im5vbmUifX19XSxbMCwxLCIiLDAseyJzdHlsZSI6eyJoZWFkIjp7Im5hbWUiOiJub25lIn19fV0sWzAsMiwiIiwyLHsic3R5bGUiOnsiaGVhZCI6eyJuYW1lIjoibm9uZSJ9fX1dLFsyLDQsIiIsMix7InN0eWxlIjp7ImhlYWQiOnsibmFtZSI6Im5vbmUifX19XSxbMCwzLCIiLDIseyJzdHlsZSI6eyJoZWFkIjp7Im5hbWUiOiJub25lIn19fV0sWzMsNCwiIiwyLHsic3R5bGUiOnsiaGVhZCI6eyJuYW1lIjoibm9uZSJ9fX1dXQ==
\[\begin{tikzcd}[ampersand replacement=\&,cramped]
	\& p \& \\
	A \& B \& C \\
	\& q
	\arrow[no head, from=1-2, to=2-1]
	\arrow[no head, from=1-2, to=2-2]
	\arrow[no head, from=1-2, to=2-3]
	\arrow[no head, from=2-1, to=3-2]
	\arrow[no head, from=2-2, to=3-2]
	\arrow[no head, from=2-3, to=3-2]
\end{tikzcd}\]
Let $U_x = X \setminus \{x\}$. Then, $U_x \cap U_y$ is always contractible, but $U_A \cap U_B \cap U_C$ is disconnected. Notice further that $\pi_1(U_x,p)=\Z$, and $\pi_1(X,p)=\Z*\Z$. Then, van Kampen would say $\Z*\Z*\Z \surj \Z*\Z$, and $\Z*\Z*\Z/\angles e \cong \Z*\Z$, but this is absurd.
\end{rmk}

We now prove van Kampen's theorem.

\begin{proof}[Proof of van Kampen]
Let $(X,p)$ be a pointed space, $X = \bigcup _{\alpha \in I}(U_\alpha,p)$, suppose $U_\alpha \cap U_\beta$ is path connected for all $\alpha, \beta$, and let $\Phi$ be the map
$$\Phi:*\pi_1(U_\alpha,p) \to \pi_1(X).$$
We claim $\Phi$ is surjective. Let $\gamma:[0,1] \to X$ be a loop at $p$, and we want to decompose it into loops in $(U_\alpha,p)$. Since $\gamma$ is continuous, the set
$$\{ \gamma^{-1}(U_\alpha)\} _{\alpha \in I}$$
is an open cover of $[0,1]$. Thus, there exists $\delta>0$ called the \href{https://en.wikipedia.org/wiki/Lebesgue%27s_number_lemma}{Lebesgue number}, such that any open subset of $[0,1]$ with diameter $<\delta$ is contained in $\gamma^{-1}(U_\alpha)$ for some $\alpha\in I$. Assume $\delta = 1/k$ for $k \in \N$, and partition the interval $[0,1]$ into $2k$ sub-intervals of length $1/2k$. Thus, each sub-interval $[\frac{i-1}{2k},\frac{i}{k}]$ is mapped via $\gamma$ into $U_\alpha$. Let $t_i = i/2k$ and $p_i = \gamma (t_i) \in U_{\alpha_i }$. Since $U_{\alpha_i}\cap U_{\alpha_{i+1}}$ is connected, there exists a curve $\mu_i$ from $p_i$ to $p$. Let $\gamma_i = \gamma|_{[t_{i-1},t_i]}$. Then,
$$\gamma = \gamma_{2k}*\cdots *\gamma_1 \sim \gamma_{2k}\cdots \gamma_{3}\mu_2^{-1}\mu_2 \gamma_2 \mu_{1}^{-1}\mu _1 \gamma_1,$$
proving surjectivity.

Now, in order to prove the desired isomorphism, it suffices to show $$\ker (\Phi) \cong  \angles{\angles{i_{\alpha \beta}(\gamma) \cdot i_{\beta\alpha}(\gamma^{-1})}}=:N.$$
Indeed, first observe that 
$$\Phi(i_{\alpha \beta}(\gamma) \cdot i_{\beta\alpha}(\gamma^{-1})) = [\gamma] [\gamma^{-1}]=e,$$
so clearly $N \subset \ker \Phi$.

For the other direction, we want to show that if $\gamma_1 *\cdots * \gamma_{k}$ is homotopic to $[e]$ as loops in $X$, then $\gamma_1 *\cdots * \gamma_k \in N$. This is provable, but the proof was confusing in lecture... I didn't feel like finishing it.
\end{proof}

\begin{exercise}
    Show that $\R^2 \setminus \{p\}$ deformation retracts to $S^1$, and conclude that $\R^2 \not \cong \R^k$ for $k>2$.
\end{exercise}

\begin{exercise}
    Compute $\pi_1(\R^2 \setminus \{p_1,\dots,p_k\})$ using van Kampen's theorem.
\end{exercise}

\begin{exercise}
    How do you hang a picture using $1$ rope and $k$ nails such that removing any nail will make it fall.
\end{exercise}

\begin{ex}[The Fundamental Group of $\R^3 \setminus S^1$]
    We can imagine the sphere of which $S^1$ is a cross section ($S_2$), and deformation retract everything outside that sphere down to it (that is, we can ignore everything fair outside our circle). We claim that $\bar B \setminus S^1$ deformation retracts to $S^2\cup L$, where $B$ is the closed ball corresponding to our sphere $S_2$, and $L$ is a line passing from one end of $S^1$ to another.

    Then, $L$ can be homotoped within the interior of $S^2$ to a closed loop touching a point of $S^2$. That is, our space retracts to $S^2 \lor S^1$, and by van Kampen's theorem
    $$\pi_1(\R \setminus \{0\}) = \pi_1(S^2 \lor S^1) = \pi_1(S^1)=\Z.$$
\end{ex}

\begin{ex}
    Now, consider $\R^3$, with multiple separate circles removed. By shrinking $\R^3$ to the balls surrounding each circle, and using the same strategy as the previous example, we see that this space is homotopy equivalent to $S^2 \lor S^1 \lor S^2 \lor S^1 \lor \cdots \lor S^2 \lor S^1$, and so the fundamental group is $\Z *\Z *\cdots * \Z$.
\end{ex}

\begin{ex}
    Now, consider $\R^3$, with two linked circles removed. We claim this is homotopy equivalent to $T^2 \lor S^2$, so the fundamental group is 
    $$\pi_1(T^2)*\pi_1(S^2)=\Z^2 *\{e\}=\Z^2.$$
The idea is as follows: we thicken one of our circles so that it becomes a torus, and we surround our torus and circle by a sphere so that the torus touches the boundary of the sphere. Then, the second circle is ``jump-roped'' around the torus so that it reaches the boundary of the sphere at that same of the intersection point.
\end{ex}

\begin{ex}
    Consider the space $S^n \setminus \{p\}$, for $n \geq 0$. We claim the fundamental group is trivial. Indeed, we may pass $S^{n-1}$ through our point, which splits $S^n$ into an ``upper-hemisphere'' $U$ and ``lower-hemisphere'' $L$. We widen $U$ and $L$ mildly so that they have non-trivial intersection, and $U \cap L \sim S^1$. Then,
    \[
    \pi_1(S^n,p)=\pi_1(U) *_{\pi_1(S^{n-1},p)}\pi_1(L),
    \]
    which by induction is trivial (since $\pi_1(S^1\setminus \{p\})=\{0\}$).
\end{ex}

\begin{prop}
Suppose $Y$ is obtained from a path-connected pointed space $(X,x_0)$ by attaching a $2$-cell, i.e.
\[
Y= X\sqcup \bar D^2/ (\varphi(x) \sim x),
\]
where $X \in \partial \bar D^2$ and $\varphi:\partial D^2 \to X$.  Let $p \in Y\cap X$ be a point such that there exists $\tilde p \in D^2$ with $\varphi(\tilde p)=p$. Let $\mu_p$ be a path in $X$ between $x_0$ and $p$. Then,
\[
\pi_1(Y,x_0) = \pi_1(X,x_0)/ \aangles{\mu_p^{-1}*\varphi * \mu_p}.
\]
\end{prop}
\begin{proof}
    Let $p'$ be a point in the interior of $D$ obtained by moving radially inward from $\tilde p$. Let $\tilde \varphi$ be a loop at $p'$, let $\varphi_1$ be a loop in $D^2$ with $\tilde \varphi$ in the interior, and $\varphi_2$ be a loop with $\tilde \varphi$ in the exterior. Let $G$ be the region within the loop $\varphi_1$, and $P$ the region outside the loop $\varphi_2$. Then, $p'$ is contained in $G \cap P$, and we may write $Y=G \cup P$. Moreover, $G \cap P$ is an annulus which is homotopy equivalent to $S^1$. It follows that
    \[
    \pi_1(Y,p') = \pi_1(G,p')*_{\pi_1(S^1,p')}\pi_1(P,p').
    \]
    Moreover, $\tilde \varphi$ is homotopic to $\bar{p'p^{-1}}\varphi \bar{p'p}$, so
    \[
    \pi_1(Y,p')\cong \pi_1(P,p')/\angles{\bar{p'p^{-1}}\varphi \bar{p'p}}.
    \]
    But, note that $P$ deformation retracts via $\mu_p\bar{p'p^{-1}}$ to give the desired result.
\end{proof}

\begin{que}
    What happens if we attach a $k$-cell for $k>2$?
\end{que}

\begin{claim}
If $Y=X\coprod D^k/\varphi(x)\sim x$ where $x\in \partial D^k$ for $k>2$, then $\pi_{1}(Y)=\pi_{1}(X)$.
That is, ``attaching $k$-cells with $k>2$ does not change the fundamental group'' 
\end{claim}

\begin{proof}
    Doing the same thing we did last time, by van Kampen's we have that 
    $$\pi_{1}(Y)=\pi_{1}(X)*_{\pi_{1}(S^{k-1})} \pi_{1}(D^k)=\pi_{1}(X)$$ 
    as both $\pi_{1}(S^{k-1})$ and $\pi_{1}(D^k)$ are trivial.
\end{proof}

Let $X$ be a finite path connected CW-complex. 

$X^0$ is just points. $X^1$ is a connected graph which is homotopic to $\bigvee_{i=1}^j S_{i}^1$ for some $j$, hence $\pi_{1}(X^1)=\mathbb{Z}^{*j}$.

$X^2$ is then obtained by attaching $2$-cells via
\[
X^2=X_{1}\coprod \coprod_{k=1}^{j_{\alpha}}\overline{D}^2_{k} \pmod{x\sim \varphi^1_{j}(x)}.
\]
Choosing $x_{0}\in X^1,$ points $p_{1},\cdots, p_{j_{\alpha}}\in \partial D^2_{k}$ and curves $\mu_{j}$ connecting $x_{0}$ to $p_{j}$ we have that 
$$\pi_{1}(X^2,x_{0})=\pi_{1}(X^1,x_{0})/ \aangles{ \mu_{i}^{-1}\varphi^1_{j}\mu_{i}}.$$
Now $\pi_{1}(X)=\pi_{1}(X^{2})$.

\begin{rmk}
    This is true without assuming finiteness. 
\end{rmk}

\begin{exercise}[Important!]
    Use van Kampen's theorem to compute
    \begin{itemize}
        \item $\pi_{1}(\mathbb{T}^2)$,
        \item $\pi_{1}(\text{Klein Bottle})$,
        \item $\pi_{1}(\mathbb{R}\mathbb{P}^2)$,
        \item $\pi_{1}(\mathbb{R}\mathbb{P}^k)$ for all $k$.
    \end{itemize}
\end{exercise}

\section{Covering Spaces}

\begin{defn}
    Let $X$ be path connected. A space $\tilde{X}$ with a continuous map $p:\tilde{X}\to X$ is called a \textbf{covering space} for $X$ if for all $x\in X$ there exists an open set $U\ni x$ such that $p^{-1}(U)=\coprod_{\alpha\in I_{x}}U_{\alpha}$ for some index set $I_{x}$, where $p:U_{\alpha}\to U$ is a homeomorphism for all $\alpha \in I_x$.
\end{defn}

\begin{ex}
    $X=S^1$, $\tilde{X} =\mathbb{R}$, $p(\theta)=e^{i\theta}$.
\end{ex}

\begin{ex}
    $X=S^1$, $\tilde{X} = S^1$, by $z\mapsto z^k$ or equivalently $\theta\to k\theta$. 
\end{ex}


\begin{ex}
    $\mathbb{R}^2\to \mathbb{R}^2/\mathbb{Z}^2\cong\mathbb{T}^2$ given by the natural projection.
\end{ex}

\begin{ex}
    $\mathbb{R}^2\to$ the Klein bottle given by the natural projection. "Think about it."
\end{ex}

\begin{ex}
We get an open cover of $S^{1}\lor S^{2}$ by turning $S^1$ into a interval. 
   Then we find another covering by "going around $A$ at twice the speed."
   We of course want another covering, so we take the first covering and turn the remaining circles into intervals as well. This last covering is connected. 
\end{ex}

\begin{prop}[Homotopy Lifting Property][label=prop:homoliftingprop]
    Suppose $\tilde{X}\to X$ is a covering map and $Y$ is a space, $f_{t}:Y\times [0,1] \to X$ is a homotopy (which inputs like $f_{.5}(y)$), and $\tilde{f}_{0}:Y\to \tilde{X}$ is such that the following diagram commutes:
% https://q.uiver.app/#q=WzAsMyxbMCwyLCJZIl0sWzIsMiwiWCJdLFsyLDAsIlxcdGlsZGV7WH0iXSxbMiwxLCJwIl0sWzAsMSwiZl8wIiwyXSxbMCwyLCJcXHRpbGRle2Z9XzAiXV0=
\[\begin{tikzcd}[ampersand replacement=\&,cramped]
	\&\& {\tilde{X}} \\
	\\
	Y \&\& X
	\arrow["p", from=1-3, to=3-3]
	\arrow["{\tilde{f}_0}", from=3-1, to=1-3]
	\arrow["{f_0}"', from=3-1, to=3-3]
\end{tikzcd}\]
Then, there exists a unique homotopy $\tilde{f_{t}}$ such that 
% https://q.uiver.app/#q=WzAsMyxbMCwyLCJZIl0sWzIsMiwiWCJdLFsyLDAsIlxcdGlsZGV7WH0iXSxbMiwxLCJwIl0sWzAsMSwiZl90IiwyXSxbMCwyLCJcXHRpbGRle2Z9X3QiXV0=
\[\begin{tikzcd}[ampersand replacement=\&,cramped]
	\&\& {\tilde{X}} \\
	\\
	Y \&\& X
	\arrow["p", from=1-3, to=3-3]
	\arrow["{\tilde{f}_t}", from=3-1, to=1-3]
	\arrow["{f_t}"', from=3-1, to=3-3]
\end{tikzcd}\]
commutes for all $t$. 
\end{prop}

\begin{proof}
Fix a point $y\in Y$, then $\operatorname{Im} (\overbrace{ t\mapsto f_{t}(y) }^{ \gamma_{y}(t) })$ is compact and for each $t$ we can choose $U_{f_{t}(y)}\subset X$ as in the definition of the covering space. (a neighborhood whose preimage under $p$ is a disjoint union of itself). Now, $\left\{ \gamma_{y}^{-1}(U_{f_{t}(y)}) \right\}_{t\in [0,1]}$ is a open cover of $[0,1]$, which is compact so it has a Lebesgue number. Then, we want to construct $\tilde{\gamma}_{y}(t):=\tilde{f_{t}}(y)$, and to do this we can trace along the interval and watch where the path ``must lift''. 

\begin{exercise}
    Prove $\tilde{f_t}$ is continuous.
\end{exercise}
\end{proof}

\begin{prop}[][label=prop:lifttoconstant]
    Let $x_0\in X$ be any point. Then, $p^{-1}(x_0)$ is discrete, so in particular any path lifting the constant path $\gamma\equiv x_0$ is constant.
\end{prop}

\begin{proof}
    Trivial from the definition, but also follows by the unique homotopy lifting property (\Cref{prop:homoliftingprop}).
\end{proof}

\begin{prop}[][label=prop:liftingtoloop]
    Suppose
    \[
    (\tilde X, \tilde{x_0}) \arr[p] (X,x_0)
    \]
    is a covering, with $X$ path connected, and fix $[\gamma] \in \pi_1(X,x_0)$. Then, $[\gamma] \in p_*\paren{\pi_1(\tilde X,\tilde{x_0})}$ if and only if $\gamma$ lifts to a loop at $\tilde{x_0}$.
\end{prop}

\begin{proof}
    If $\gamma$ lifts to a loop $\tilde \gamma$ in $\pi_1(\tilde{X},\tilde{x_0})$, then $\gamma = p \circ \tilde \gamma$, hence $p_*\paren{[\tilde \gamma]} = \gamma$. On the other hand, if $[\gamma] \in p_*\paren{\pi_1(\tilde X,\tilde{x_0})}$, then $\gamma$ is homotopic to $p \circ \tilde \mu$ for some loop $\tilde \mu$ at $\tilde{x_0}$ in $\tilde X$. Denoting by $f_t$ this homotopy such that $f_0 = p \circ \tilde \mu$ and $f_1 = \gamma$, by the homotopy lifting property, there exists a unique $\tilde{f_t}$ lifting $f_t$ such that $\tilde{f_0}=\tilde{\gamma}$.

% https://q.uiver.app/#q=WzAsNixbMiwwLCJcXHRpbGRlIFgiXSxbMCwyLCJbMCwxXSJdLFsyLDIsIlgiXSxbNCwyLCJbMCwxXSJdLFs2LDAsIlxcdGlsZGUgWCJdLFs2LDIsIlgiXSxbMSwwLCJcXHRpbGRle1xcbXV9Il0sWzAsMiwicCJdLFsxLDIsImZfMD1wIFxcY2lyYyBcXHRpbGRlIFxcbXUiLDJdLFszLDQsIlxcdGlsZGUge2ZfdH0iXSxbNCw1LCJwIl0sWzMsNSwiZl90IiwyXV0=
\[\begin{tikzcd}[ampersand replacement=\&,cramped]
	\&\& {\tilde X} \&\&\&\& {\tilde X} \\
	\\
	{[0,1]} \&\& X \&\& {[0,1]} \&\& X
	\arrow["p", from=1-3, to=3-3]
	\arrow["p", from=1-7, to=3-7]
	\arrow["{\tilde{\mu}}", from=3-1, to=1-3]
	\arrow["{f_0=p \circ \tilde \mu}"', from=3-1, to=3-3]
	\arrow["{\tilde {f_t}}", from=3-5, to=1-7]
	\arrow["{f_t}"', from=3-5, to=3-7]
\end{tikzcd}\]

    We claim $\tilde{f_1}$ is a loop at $\tilde{x_0}$ lifting $\gamma$. Indeed, as $f_t$ is base-point preserving, $t \mapsto f_t(1)$ is a constant map. Hence, by \Cref{prop:lifttoconstant}, $f_t$ lifts to a constant path with
    \[
    \tilde{f_1}(1)=\tilde{f_0}(1)=\tilde{\mu}(1)=\tilde{x_0}.
    \]
\end{proof}

\begin{que}
When do two elements $[\gamma_1],[\gamma_2] \in \pi_1(X,x_0)$ satisfy $\tilde{\gamma_1}(1) = \tilde{\gamma_2}(1)$ in $p^{-1}(x_0)$, assuming $\tilde{\gamma_1}(0)=\tilde{\gamma_2}(0)=\tilde{x_0}$?
\end{que}

Well, knowing $\tilde{\gamma_1}(1) = \tilde{\gamma_2}(1)$, we may apply \Cref{prop:liftingtoloop}, hence $\tilde{\gamma_2}^{-1}*\tilde{\gamma_1}$ is a loop at $\tilde{x_0}$ if and only if
\[
\braks{\gamma_2^{-1}*\gamma_1}\in p_*\paren{\pi_1(X,x_0)} \iff \gamma_1 \in \braks{\gamma_2}p_*\paren{\pi_1(\tilde X,\tilde{x_0})},
\] 
i.e. when we have equivalence in the quotient space.

\begin{cor}
    $p^{-1}(x_0)$ is in bijective correspondence with
    \[
    \pi_1(X,x_0)/p_*\paren{\pi_1(\tilde X,\tilde{x_0})}.
    \]
\end{cor}

\begin{que}
When can $f$ be lifted to a map
\[
\tilde f:(Y,y_0)\to(\tilde X, \tilde{x_0})?
\]
That is, when can me make the following diagram commute?:
% https://q.uiver.app/#q=WzAsMyxbMCwyLCIoWSx5XzApIl0sWzIsMCwiKFxcdGlsZGUgWCwgXFx0aWxkZXt4XzB9KSJdLFsyLDIsIihYLHhfMCkiXSxbMCwxLCJcXHRpbGRlIGYiLDAseyJzdHlsZSI6eyJib2R5Ijp7Im5hbWUiOiJkYXNoZWQifX19XSxbMCwyLCJmIiwyXSxbMSwyLCJwIl1d
\[\begin{tikzcd}[ampersand replacement=\&,cramped]
	\&\& {(\tilde X, \tilde{x_0})} \\
	\\
	{(Y,y_0)} \&\& {(X,x_0)}
	\arrow["p", from=1-3, to=3-3]
	\arrow["{\tilde f}", dashed, from=3-1, to=1-3]
	\arrow["f"', from=3-1, to=3-3]
\end{tikzcd}\]
\end{que}

\begin{thm}
    Suppose $Y$ is path connected and locally path connected. Then, a map $f:(Y,y_0) \to (X,x_0)$ can be lifted to $\tilde f:(Y,y_0) \to (\tilde X,\tilde {x_0})$ if and only if
    \[
    f_*\paren{\pi_1(Y,y_0)} \subset p_*\paren{\pi_1(\tilde X,\tilde {x_0})}
    \]
\end{thm}


\begin{proof}
    If a lift exists, i.e. there exists $\tilde f$ such that $p \circ \tilde f = f$, then
    \[
    f_* =p_* \circ \tilde{f_*} \implies \Img(f_*) \subset \Img (p_*),
    \]
    proving the first direction. For the other direction, fix $y \in Y$, and take a path $\gamma$ connecting $y_0$ to $y$. Then, $f \circ \gamma$ is a curve in $X$ starting at $x_0$, so by the homotopy lifting property there exists $\mu_\gamma(t)$ lifting $f(\gamma(t))$. Define $\tilde f(y):=\mu_\gamma(1)$. We now need to show that this $\tilde f$ is independent of the choice of $\gamma$. Continued next time.
\end{proof}


\end{document}
